\documentclass[main.tex]{subfiles}

\begin{document}

\section{Related Work}
\label{sec:related_work}

\subsection{Brain Tumor MRI Classification and Transfer Learning}

Brain tumor MRI classification is commonly addressed as a supervised medical-image classification problem in which an imaging study, slice, or region is assigned to a diagnostic or tumor-related category. Because labeled medical-image datasets are often limited relative to the capacity of modern visual architectures, practical workflows frequently rely on transfer learning, fine-tuning, or frozen-feature reuse rather than learning the full visual representation from scratch \cite{dorfner2025review,yosinski2014transferable,bar2024frofa}. In this setting, a trained backbone may be used directly through its native softmax head, or indirectly as a feature extractor whose embeddings are passed to a downstream classifier. The latter setting is especially relevant for post-hoc evaluation, because it defines the final-embedding baseline: if the deepest representation is already sufficient, a multilayer method must demonstrate value beyond that simpler alternative.

\subsection{Embedding-Space Classifiers and Distance-Based Reuse}

An embedding is a vector representation intended to preserve task-relevant information about an input image. In feature-reuse workflows, embeddings can be passed to downstream classifiers, retrieval procedures, calibration modules, or adaptation methods \cite{bar2024frofa,li2025embeddings}. Their downstream behavior depends on the dimensionality of the vector space, the scale of its coordinates, and the geometry it induces among samples, that is, which images become close or distant in the resulting feature space.

This geometric view motivates the use of distance-based classifiers. The k-nearest neighbors classifier (kNN) is a transparent downstream method because prediction depends directly on the labels of nearby training samples \cite{cover1967nearest}. Metric-learning approaches, such as Neighborhood Components Analysis (NCA) and Large-Margin Nearest Neighbor learning (LMNN), also emphasize this relation between embedding geometry and neighbor-based prediction by explicitly adapting a representation or metric for local classification \cite{goldberger2004nca,weinberger2009lmnn}. Other classical classifiers, including support vector machines, random forests, and gradient-boosted trees, provide additional final-embedding baselines once a feature vector is available \cite{cortes1995svm,breiman2001random,chen2016xgboost}.

Distance-based reuse is sensitive to scale, dimensionality, and redundant coordinates. In high-dimensional spaces, pairwise distances may become less discriminative, a phenomenon commonly associated with distance concentration \cite{beyer1997nearest}. This limitation is directly relevant to multilayer CNN embeddings: adding internal representations may introduce complementary information, redundant feature blocks, heterogeneous scales, or poorly aligned distance structures. For this reason, multilayer embedding reuse requires a selection or fusion rule controlled by a valid evaluation procedure.

\subsection{Internal CNN Representations and Multilayer Feature Fusion}

The final embedding of a trained convolutional neural network (CNN) is a natural descriptor because it is directly connected to the original classification head. CNNs also expose multiple internal representations during the same forward pass, and prior work on feature visualization and transferability has shown that these representations vary across network depth \cite{zeiler2014visualizing,yosinski2014transferable}. This depth-dependent behavior motivates the use of internal representations as candidate embeddings rather than treating the final vector as the only possible descriptor.

Several works have explored representations from more than one network depth. Fradi et al. extracted pooled convolutional features from multiple CNN layers, normalized and concatenated them, and then applied dimensionality reduction and feature selection for texture classification \cite{fradi2021multilayer}. Gomes et al. proposed a multilayer feature-fusion method in a few-shot setting, where classification depends on comparing examples in an embedding space \cite{gomes2023fmlf}. Head2Toe uses intermediate representations for transfer learning by selecting feature groups from multiple depths \cite{evci2022head2toe}. In medical-image classification, patch-level, hierarchical, multiscale, and attention-based methods further support the idea that information from different spatial or semantic levels can improve classification when the fusion mechanism is aligned with the task \cite{muezzinoglu2023patchresnet,ullah2025hierarchical,sankari2025hierarchical,he2025hemf,zhang2025context}.

These works motivate internal-representation reuse and show that multilayer fusion can be implemented at different stages of the pipeline. Architectural fusion combines features during model training by adding trainable modules, attention blocks, or multiscale branches. Decision-level fusion combines predictions after separate classifiers have been applied. Representation-level fusion constructs a shared descriptor before classification, often through concatenation or projection. Distance-level fusion provides a related but distinct formulation: each layer remains a separate metric view, and the distances induced by selected views are combined before applying a distance-based decision rule. The present work focuses on this last setting.

\subsection{Multiview, Kernel, and Projection-Based Fusion}

Multilayer CNN embeddings can also be interpreted as multiple paired views of the same image, where each view corresponds to a different layer representation. Under this interpretation, multiview representation learning provides relevant reference methods. Generalized canonical correlation analysis (GCCA) seeks a common representation across multiple sets of variables \cite{carroll1968gca,kettenring1971canonical}. Generalized multiview analysis (GMA) and multiview discriminant analysis (MvDA) extend this family by incorporating class-discriminative objectives into multiview projection \cite{sharma2012gma,kan2016mvda}. Multiple kernel learning provides another route by combining kernels defined over different views; EasyMKL is a scalable example of this family \cite{aiolli2015easymkl}.

These methods are relevant because they combine several views through projection or kernel weighting. Projection-based multiview methods learn a shared latent representation; kernel-based methods combine view-specific similarity functions. The proposed method follows a different route: it selects a compact set of layer-wise Euclidean distance spaces using validation evidence and performs distance-weighted neighbor voting after freezing the selected configuration. In the experimental comparison, these methods are treated as principled stage-level reference implementations over the same saved CNN embedding views.

\subsection{Post-Hoc Metric and Prototype Reuse}

Post-hoc metric and prototype methods provide complementary forms of representation reuse. NCA followed by kNN tests whether a supervised metric projection over saved embeddings improves local neighborhood classification \cite{goldberger2004nca}. Prototype-based methods classify samples by comparing them with class representatives instead of relying directly on all training examples. KMEx fits class-wise prototypes on a trained model's embedding without retraining the model \cite{gautam2024kmex}, while recent spatial prototype approaches assign image regions to prototypes to provide post-hoc self-explanations \cite{boubekki2026posthoc}.

These methods are relevant because they also reuse information extracted from an existing network. Their primary mechanisms differ from multilayer distance fusion. Metric-projection methods modify or project the representation before classification, whereas prototype methods rely on representative class anchors or prototype assignments. The proposed method keeps the selected layer spaces as metric views, combines their layer-wise distances, and applies frozen distance-weighted kNN inference.

\subsection{Positioning of \method{}}

The preceding literature motivates a post-hoc multilayer classifier evaluated against the native softmax head, final-embedding classifiers, multilayer and multiview reference methods, and metric or prototype-based alternatives under the same train/validation/test protocol. Such a comparison separates possible sources of improvement, including generic view combination, improved embedding geometry, useful intermediate depth, and model-selection effects. \fullmethod{} (\method{}) is positioned within this space as a validation-controlled distance-fusion approach: it reuses internal representations from a trained CNN while keeping the backbone architecture fixed, preserving the trained feature extractor during the post-hoc stage, and reserving the test set for final evaluation.

\end{document}
